% archon:covers AlgebraicJacobian/RiemannRoch/CurveKrullDim.lean

\chapter{The Krull dimension of a smooth curve is at most one}
\label{chap:RiemannRoch_CurveKrullDim}

\providecommand{\krullDim}{\dim}
\providecommand{\coheight}{\operatorname{coheight}}
\providecommand{\height}{\operatorname{height}}
\providecommand{\trdeg}{\operatorname{trdeg}}
\providecommand{\Frac}{\operatorname{Frac}}

% SOURCE: [Stacks Project] dimension theory, read from
% references/stacks-regular-dvr-algebra.tex (a verbatim copy of the Stacks
% `algebra.tex').  The crux \cref{lem:chart_ringKrullDim_le_one} uses the
% embedding-dimension inequality $\dim R \le \dim_\kappa(\mathfrak m/\mathfrak m^2)$
% at each maximal ideal; the residue field is $\bar k$ by the Nullstellensatz, and
% the cotangent dimension bound is reused from \cref{chap:RiemannRoch_OcOfD}.
% The Cohen--Seidenberg keystone \cref{lem:ringKrullDim_le_of_moduleFinite_injective}
% (Tag 00OJ, `lemma-integral-dim-up', l.27352) is retained as reusable infrastructure.

\section{Purpose and standing hypotheses}

This chapter supplies the single dimension-theoretic input required by the
codimension-one regularity witness of \cref{chap:RiemannRoch_OcOfD}: the bound
\[
  \krullDim C \;\le\; 1
\]
on the Krull dimension of the underlying space of a smooth, relative-dimension-one
curve \(C\) over an algebraically closed field \(\bar k\).  Here \(\krullDim\) is the
order-theoretic Krull dimension of the topological space \(\lvert C\rvert\) under its
specialisation order.  The bound is consumed --- as the hypothesis \(\krullDim C \le
1\) of the cotangent pin
\cref{lem:cotangentSpace_finrank_eq_one_of_smooth} --- to license the passage
``coheight one \(\Rightarrow\) closed point'' that underlies the discrete-valuation
property at each codimension-one stalk.

\paragraph{Standing hypotheses.}  Throughout, \(\bar k\) is an algebraically closed
field and \(C : \Over(\Spec \bar k)\) is a curve over \(\bar k\) whose structure
morphism \(C.\mathrm{hom} : C \to \Spec \bar k\) is smooth of relative dimension one
(\(\texttt{SmoothOfRelativeDimension}\ 1\ C.\mathrm{hom}\)) and locally of finite type
(\(\texttt{LocallyOfFiniteType}\ C.\mathrm{hom}\)), and whose underlying scheme
\(C.\mathrm{left}\) is integral (\(\texttt{IsIntegral}\ C.\mathrm{left}\)).

\paragraph{Strategy.}  The argument has three layers.

\begin{enumerate}
  \item \emph{Outer reduction (free).}  Krull dimension of a space is the supremum of
    the coheights of its points, and for a scheme the coheight of a point equals the
    Krull dimension of its local ring.  This reduces \(\krullDim C \le 1\) to a
    \emph{pointwise} bound \(\krullDim \mathcal O_{C,z} \le 1\) on every stalk.
  \item \emph{Localisation push (free).}  At each point \(z\) the stalk
    \(\mathcal O_{C,z}\) is the localisation of a standard-smooth affine chart ring
    \(S\) at a prime \(\mathfrak p_z\); the Krull dimension of a localisation at a
    prime is the height of that prime, which is at most the Krull dimension of \(S\).
    This reduces the pointwise bound to \(\krullDim S \le 1\) for the chart ring.
  \item \emph{The crux (embedding-dimension route).}  For a
    standard-smooth, relative-dimension-one \(\bar k\)-algebra \(S\) which is an
    integral domain, \(\krullDim S \le 1\).  Write the Krull dimension as the
    supremum of the prime heights of the \emph{maximal} ideals.  At each maximal
    ideal \(\mathfrak m\), the localisation \(S_\mathfrak m\) is a Noetherian local
    ring whose residue field is \(\bar k\) --- because \(S\) is finite type over the
    algebraically closed field \(\bar k\), the Nullstellensatz forces every maximal
    quotient to equal \(\bar k\), with \emph{no} appeal to the global dimension bound.
    Krull's height theorem bounds \(\krullDim S_\mathfrak m\) by the embedding
    dimension \(\dim_{\bar k}(\mathfrak m S_\mathfrak m / \mathfrak m^2 S_\mathfrak
    m)\), and the conormal bridge of \cref{chap:RiemannRoch_OcOfD} (already
    proved, reused here at the maximal ideals of the chart) shows this cotangent
    dimension is one.  Hence every maximal height is \(\le 1\) and \(\krullDim S \le
    1\).
\end{enumerate}

Every external step invoked in the crux is present in Mathlib, and the
one-dimensionality of the cotangent space is the conormal corollary already
established for the DVR-stalk substrate.  The Cohen--Seidenberg keystone
\cref{lem:ringKrullDim_le_of_moduleFinite_injective} and its incomparability engine
\cref{lem:comap_lt_comap_of_isIntegral} are retained below as general reusable
infrastructure.

\section{Mathlib dependency anchors}

The outer reduction and the localisation push are assembled entirely from Mathlib
results, recorded here as dependency anchors.

\begin{lemma}[Krull dimension as a supremum of coheights]
  \label{lem:krullDim_eq_iSup_coheight_mathlib}
  \lean{Order.krullDim_eq_iSup_coheight}
  \mathlibok
  \textit{Provided by Mathlib (\texttt{Mathlib.Order.KrullDimension}).}
  For a preorder \(\alpha\), the Krull dimension of \(\alpha\) is the supremum over
  all \(a : \alpha\) of the coheights:
  \(\krullDim \alpha = \bigsqcup_{a} \coheight(a)\) in \(\mathrm{WithBot}\,\mathbb
  N_\infty\).
\end{lemma}

\begin{lemma}[Stalk Krull dimension as a localisation height]
  \label{lem:ringKrullDim_eq_height_mathlib}
  \lean{IsLocalization.AtPrime.ringKrullDim_eq_height}
  \mathlibok
  \textit{Provided by Mathlib (\texttt{Mathlib.RingTheory.Ideal.Height}).}
  Let \(R\) be a commutative ring and \(I \subseteq R\) a prime ideal.  If \(A\) is an
  \(R\)-algebra that is a localisation of \(R\) at \(I\)
  (\(\texttt{IsLocalization.AtPrime}\ A\ I\)), then \(\krullDim A = \height I\).
\end{lemma}

\begin{lemma}[Prime height is bounded by Krull dimension]
  \label{lem:height_le_ringKrullDim_mathlib}
  \lean{Ideal.primeHeight_le_ringKrullDim}
  \mathlibok
  % NOTE: \lean{} targets `Ideal.primeHeight_le_ringKrullDim` (with `Ideal.height_eq_primeHeight`),
  % the declarations that `krullDim_curve_le_one` actually invokes.
  \textit{Provided by Mathlib (\texttt{Mathlib.RingTheory.Ideal.Height}).}
  For a commutative ring \(R\) and a prime ideal \(I \subseteq R\), the height of
  \(I\) is at most the Krull dimension of \(R\): \(\height I \le \krullDim R\).
\end{lemma}

\begin{lemma}[Krull dimension as a supremum over maximal ideals]
  \label{lem:sup_primeHeight_maximal_mathlib}
  \lean{Ideal.sup_primeHeight_of_maximal_eq_ringKrullDim}
  \mathlibok
  \textit{Provided by Mathlib (\texttt{Mathlib.RingTheory.Ideal.Height}).}
  For a nontrivial commutative ring \(R\), the Krull dimension of \(R\) is the
  supremum of the prime heights of its \emph{maximal} ideals:
  \(\krullDim R = \bigsqcup_{\mathfrak m \text{ maximal}} \height \mathfrak m\).
\end{lemma}

% SOURCE: [Stacks Project], `algebra.tex' (inline prose preceding
% \label{definition-regular-local}, lines 14604--14613):
% SOURCE QUOTE: "Let $(R, \mathfrak m)$ be a Noetherian local ring. From the above it
% is clear that $\mathfrak m$ cannot be generated by fewer than $\dim(R)$ variables.
% By Nakayama's Lemma \ref{lemma-NAK} the minimal number of generators of $\mathfrak m$
% equals $\dim_{\kappa(\mathfrak m)} \mathfrak m/\mathfrak m^2$. Hence we have the
% following fundamental inequality $\dim(R) \leq \dim_{\kappa(\mathfrak m)}
% \mathfrak m/\mathfrak m^2$."
\begin{lemma}[Krull's height theorem: dimension below embedding dimension]
  \label{lem:ringKrullDim_le_spanFinrank_mathlib}
  \lean{ringKrullDim_le_spanFinrank_maximalIdeal}
  \mathlibok
  \textit{Provided by Mathlib (\texttt{Mathlib.RingTheory.Ideal.KrullsHeightTheorem}).}
  For a Noetherian local ring \(R\), the Krull dimension of \(R\) is at most the
  minimal number of generators of its maximal ideal:
  \(\krullDim R \le \operatorname{spanFinrank}(\mathfrak m_R)\).
\end{lemma}

\begin{lemma}[Minimal generators of the maximal ideal equal the cotangent dimension]
  \label{lem:spanFinrank_eq_finrank_cotangent_mathlib}
  \lean{IsLocalRing.spanFinrank_maximalIdeal_eq_finrank_cotangentSpace}
  \mathlibok
  \textit{Provided by Mathlib (\texttt{Mathlib.Algebra.Module.SpanRankOperations}).}
  For a Noetherian local ring \(R\), the minimal number of generators of the maximal
  ideal equals the dimension of the cotangent space over the residue field:
  \(\operatorname{spanFinrank}(\mathfrak m_R) = \dim_{\kappa(R)}(\mathfrak m_R/\mathfrak
  m_R^2)\).  Combined with \cref{lem:ringKrullDim_le_spanFinrank_mathlib} this gives
  \(\krullDim R \le \dim_{\kappa(R)}(\mathfrak m_R/\mathfrak m_R^2)\).
\end{lemma}

\begin{lemma}\leanok
[Cotangent dimension one at a standard-smooth localisation]
  \label{lem:cotangent_localization_eq_one_bridge}
  \lean{Algebra.finrank_cotangentSpace_localization_eq_one_of_isStandardSmooth_dim_one}
  \textit{Established in \cref{chap:RiemannRoch_OcOfD} (reused here).}
  Let \(S\) be a relative-dimension-one standard-smooth \(\bar k\)-algebra and let
  \(T\) be a localisation of \(S\) at a prime, a local ring whose residue field is
  \(\bar k\)-rational (encoded by the two vanishings \(\Omega_{\kappa(T)/\bar k} = 0\)
  and \(H_1^{\mathrm{cot}}(\bar k, \kappa(T)) = 0\)).  Then
  \(\dim_{\kappa(T)}(\mathfrak m_T/\mathfrak m_T^2) = 1\).
\end{lemma}

\begin{lemma}[Krull dimension of a polynomial ring over a Noetherian base]
  \label{lem:mvpoly_ringKrullDim_mathlib}
  \lean{MvPolynomial.ringKrullDim_of_isNoetherianRing}
  \mathlibok
  \textit{Provided by Mathlib (\texttt{Mathlib.RingTheory.KrullDimension.Polynomial}).}
  For a Noetherian commutative ring \(R\) and a finite index type \(\iota\),
  \(\krullDim\bigl(R[\iota]\bigr) = \krullDim R + \lvert\iota\rvert\), where
  \(R[\iota]\) is the polynomial ring \(\texttt{MvPolynomial}\ \iota\ R\) and
  \(\lvert\iota\rvert\) is the cardinality of \(\iota\).  In particular, over a field
  \(\bar k\) (where \(\krullDim \bar k = 0\)) one has \(\krullDim(\bar k[\iota]) =
  \lvert\iota\rvert\).
\end{lemma}

\begin{lemma}[Noether normalisation, module-finite injective form]
  \label{lem:noether_normalization_mathlib}
  \lean{exists_finite_inj_algHom_of_fg}
  \mathlibok
  \textit{Provided by Mathlib (\texttt{Mathlib.RingTheory.NoetherNormalization}).}
  Let \(k\) be a field and \(R\) a nonzero, finitely generated \(k\)-algebra
  (\(\texttt{Algebra.FiniteType}\ k\ R\)).  Then there is a natural number \(s\) and a
  \(k\)-algebra homomorphism \(g : k[X_1,\dots,X_s] \to R\) from the polynomial ring in
  \(s\) variables that is injective and module-finite (\(g.\texttt{Finite}\)).
\end{lemma}

The project's own bridge from a scheme point to its stalk dimension is recorded next;
it is established in the file \texttt{AlgebraicJacobian/Albanese/CoheightBridge.lean}.

\begin{lemma}\leanok
[Coheight of a scheme point equals the stalk Krull dimension]
  \label{lem:ringKrullDim_stalk_eq_coheight}
  \lean{AlgebraicGeometry.Scheme.ringKrullDim_stalk_eq_coheight}
  \uses{lem:ringKrullDim_eq_height_mathlib}
  For any scheme \(X\) and any point \(z : X\), the Krull dimension of the stalk
  \(\mathcal O_{X,z}\) equals the coheight of \(z\) in the underlying space (with the
  specialisation preorder): \(\krullDim \mathcal O_{X,z} = \coheight(z)\).
\end{lemma}

\begin{proof}
  \uses{lem:ringKrullDim_eq_height_mathlib}
  \leanok
  Choose an affine open \(U = \Spec \Gamma(X,U) \ni z\) and let \(\mathfrak p\) be the
  prime of \(\Gamma(X,U)\) corresponding to \(z\).  The stalk \(\mathcal O_{X,z}\) is
  the localisation of \(\Gamma(X,U)\) at \(\mathfrak p\), so by
  \cref{lem:ringKrullDim_eq_height_mathlib} its Krull dimension is the height of
  \(\mathfrak p\) computed in \(\Spec \Gamma(X,U)\).  On the other hand the coheight
  of \(z\), being a property of the specialisation order, is unchanged on passing to
  the open subspace \(U\), and under the homeomorphism \(U \cong \Spec \Gamma(X,U)\)
  the coheight of \(z\) equals the height of \(\mathfrak p\).  Combining the two
  identifications gives \(\krullDim \mathcal O_{X,z} = \coheight(z)\).
\end{proof}

\section{The keystone: Cohen--Seidenberg upper bound}

\emph{The two lemmas of this section are general Cohen--Seidenberg infrastructure;
they are not used in the proof of \cref{lem:chart_ringKrullDim_le_one} but are
available for reuse.}

% SOURCE: [Stacks Project], incomparability for integral extensions, Tag 00H8
% (`lemma-integral-no-inclusion'); the project lemma generalises Mathlib's
% `Ideal.IntegralClosure.comap_lt_comap` (which is stated only for an
% `integralClosure R A`) to an arbitrary integral algebra.
\begin{lemma}\leanok
[Incomparability: strict contraction along an integral extension]
  \label{lem:comap_lt_comap_of_isIntegral}
  \lean{comap_lt_comap_of_isIntegral}
  \textit{Source: Stacks Project, Tag 00H8 (incomparability).}
  Let \(A\) be an integral \(R\)-algebra and let \(I \subsetneq J\) be prime ideals of
  \(A\) with \(I\) prime.  Then the contractions are strict:
  \(I \cap R \subsetneq J \cap R\) (in algebra notation,
  \(\operatorname{comap}(\mathrm{alg}) I \subsetneq \operatorname{comap}(\mathrm{alg})
  J\)).
\end{lemma}

\begin{proof}\leanok
  Monotonicity of contraction gives \(I \cap R \subseteq J \cap R\).  For strictness,
  suppose \(I \cap R = J \cap R = \mathfrak p\).  Pass to the integral extension of
  domains \(R/\mathfrak p \hookrightarrow A/I\); the image \(J/I\) is a nonzero ideal
  whose contraction to \(R/\mathfrak p\) is \(\bot\).  But a nonzero ideal of a domain
  integral over a domain contracts to a nonzero ideal (Stacks Tag 00H8 /
  \texttt{Ideal.under\_ne\_bot}), a contradiction.  Hence the contraction is strict.
\end{proof}


% SOURCE: [Stacks Project], Tag 00OJ (`lemma-integral-dim-up'), read from
% references/stacks-regular-dvr-algebra.tex, lines 27351--27362.  The going-up
% mechanism it depends on is Tag 00OH (`lemma-dimension-going-up', lines
% 27316--27336) and the equality version is Tag 00OK (`lemma-integral-sub-dim-equal',
% lines 27364--27376).
% SOURCE QUOTE: "Suppose that $R \to S$ is a ring map such that $S$ is integral over
% $R$. Then $\dim (R) \geq \dim(S)$, and every closed point of $\Spec(S)$ maps to a
% closed point of $\Spec(R)$."
\begin{lemma}\leanok
[Krull dimension does not increase along a module-finite injective extension]
  \label{lem:ringKrullDim_le_of_moduleFinite_injective}
  \lean{ringKrullDim_le_of_moduleFinite_injective}
  \textit{Source: Stacks Project, Tag 00OJ.}
  Let \(f : R \to S\) be an injective homomorphism of commutative rings such that
  \(S\) is module-finite over \(R\) (equivalently, \(f\) is finite; in particular
  \(S\) is integral over \(R\)).  Then \(\krullDim S \le \krullDim R\).
\end{lemma}

% SOURCE QUOTE PROOF: "Immediate from Lemmas \ref{lemma-integral-no-inclusion} and
% \ref{lemma-going-up-maximal-on-top} and the definitions."
% (The going-up chain mechanism this defers to is Tag 00OH, whose proof reads:
% "Assume going up. Take any chain $\mathfrak p_0 \subset \mathfrak p_1 \subset \ldots
% \subset \mathfrak p_e$ of prime ideals in $R$. By surjectivity we may choose a prime
% $\mathfrak q_0$ mapping to $\mathfrak p_0$. By going up we may extend this to a chain
% of length $e$ of primes $\mathfrak q_i$ lying over $\mathfrak p_i$. Thus $\dim(S)
% \geq \dim(R)$. The case of going down is exactly the same.")
\begin{proof}
  \uses{lem:comap_lt_comap_of_isIntegral}
\leanok
  Since \(S\) is module-finite over \(R\), it is integral over \(R\); the statement
  is therefore the upper bound of the Cohen--Seidenberg theorem (Stacks~Tag~00OJ).
  Concretely, let
  \[
    \mathfrak q_0 \subsetneq \mathfrak q_1 \subsetneq \dots \subsetneq \mathfrak q_e
  \]
  be a strict chain of primes in \(S\), and let \(\mathfrak p_i = f^{-1}(\mathfrak
  q_i)\) be its contraction to \(R\).  The chain of contractions is increasing; it is
  moreover \emph{strict}, because an integral extension satisfies incomparability ---
  two distinct primes of \(S\) with the same contraction to \(R\) would be comparable
  and equal, contradicting \(\mathfrak q_i \subsetneq \mathfrak q_{i+1}\).  Hence
  \(\mathfrak p_0 \subsetneq \dots \subsetneq \mathfrak p_e\) is a strict chain of the
  same length \(e\) in \(R\), so every chain length realised in \(S\) is realised in
  \(R\), giving \(\krullDim S \le \krullDim R\).  Injectivity of \(f\) is what makes
  the contraction map \(\Spec S \to \Spec R\) dominant, ensuring no length is lost; it
  is the going-up direction (Stacks~Tag~00OH), together with this incomparability
  bound, that yields the equality \(\krullDim S = \krullDim R\) of Stacks~Tag~00OK,
  of which only the \(\le\) half is needed here.
\end{proof}

\section{The crux: dimension of a standard-smooth chart (embedding-dimension route)}

\begin{lemma}\leanok
[Krull dimension of a relative-dimension-one standard-smooth chart]
  \label{lem:chart_ringKrullDim_le_one}
  \lean{AlgebraicGeometry.chart_ringKrullDim_le_one}
  \uses{lem:sup_primeHeight_maximal_mathlib, lem:ringKrullDim_le_spanFinrank_mathlib,
        lem:spanFinrank_eq_finrank_cotangent_mathlib, lem:ringKrullDim_eq_height_mathlib,
        lem:cotangent_localization_eq_one_bridge}
  Let \(S\) be a standard-smooth \(\bar k\)-algebra of relative dimension one over an
  algebraically closed field \(\bar k\), which is an integral domain.  Then
  \(\krullDim S \le 1\).
\end{lemma}

% SOURCE: [Stacks Project], Hilbert Nullstellensatz, `algebra.tex'
% (\label{theorem-nullstellensatz}):
% SOURCE QUOTE: "Let $k$ be a field. (1) For any maximal ideal $\mathfrak m \subset
% k[x_1, \ldots, x_n]$ the field extension $\kappa(\mathfrak m)/k$ is finite. [...]
% The same is true in any finite type $k$-algebra."
% SOURCE QUOTE PROOF: "It is enough to prove part (1) of the theorem for the case of a
% polynomial algebra $k[x_1, \ldots, x_n]$, because any finitely generated $k$-algebra
% is a quotient of such a polynomial algebra."
\begin{proof}
  \uses{lem:sup_primeHeight_maximal_mathlib, lem:ringKrullDim_le_spanFinrank_mathlib,
        lem:spanFinrank_eq_finrank_cotangent_mathlib, lem:ringKrullDim_eq_height_mathlib,
        lem:cotangent_localization_eq_one_bridge}
        \leanok
  Since a standard-smooth algebra is of finite presentation, \(S\) is of finite type
  over \(\bar k\), and being a domain it is nontrivial.  By
  \cref{lem:sup_primeHeight_maximal_mathlib},
  \[
    \krullDim S \;=\; \bigsqcup_{\mathfrak m \text{ maximal}} \height \mathfrak m ,
  \]
  so it suffices to prove \(\height \mathfrak m \le 1\) for every maximal ideal
  \(\mathfrak m \subseteq S\).

  Fix a maximal ideal \(\mathfrak m\) and pass to the localisation \(S_\mathfrak m\),
  a Noetherian local domain.  By \cref{lem:ringKrullDim_eq_height_mathlib} the height
  of \(\mathfrak m\) is the Krull dimension of \(S_\mathfrak m\).  Krull's height
  theorem \cref{lem:ringKrullDim_le_spanFinrank_mathlib} together with
  \cref{lem:spanFinrank_eq_finrank_cotangent_mathlib} bounds this by the cotangent
  dimension:
  \[
    \height \mathfrak m \;=\; \krullDim S_\mathfrak m
      \;\le\; \operatorname{spanFinrank}(\mathfrak m S_\mathfrak m)
      \;=\; \dim_{\kappa(S_\mathfrak m)}\bigl(\mathfrak m S_\mathfrak m / \mathfrak m^2
      S_\mathfrak m\bigr).
  \]

  It remains to show this cotangent dimension is \(1\).  Because \(S\) is finite type
  over the algebraically closed field \(\bar k\) and \(\mathfrak m\) is maximal, the
  Nullstellensatz gives a residue-field isomorphism \(\kappa(S_\mathfrak m) = S/\mathfrak
  m \cong \bar k\); concretely the composite \(S \to S_\mathfrak m \to \kappa(S_\mathfrak
  m)\) is surjective, so \(\kappa(S_\mathfrak m)\) is finite type over \(\bar k\), hence
  equal to \(\bar k\) (algebraically closed).  Rationality of the residue field supplies
  the two vanishings \(\Omega_{\kappa(S_\mathfrak m)/\bar k} = 0\) and
  \(H_1^{\mathrm{cot}}(\bar k, \kappa(S_\mathfrak m)) = 0\), which are exactly the
  hypotheses of the conormal bridge
  \cref{lem:cotangent_localization_eq_one_bridge}.  Applying that bridge to the
  localisation \(S_\mathfrak m\) of the relative-dimension-one standard-smooth algebra
  \(S\) gives \(\dim_{\kappa(S_\mathfrak m)}(\mathfrak m S_\mathfrak m/\mathfrak
  m^2 S_\mathfrak m) = 1\).

  Therefore \(\height \mathfrak m \le 1\) for every maximal \(\mathfrak m\), and taking
  the supremum yields \(\krullDim S \le 1\).
\end{proof}

\section{The dimension bound for a smooth curve}

\begin{theorem}\leanok
[Krull dimension of a smooth curve is at most one]
  \label{thm:krullDim_curve_le_one}
  \lean{AlgebraicGeometry.krullDim_curve_le_one}
  \uses{lem:krullDim_eq_iSup_coheight_mathlib, lem:ringKrullDim_stalk_eq_coheight,
        lem:stalk_standardSmooth_localization_of_smooth,
        lem:ringKrullDim_eq_height_mathlib, lem:height_le_ringKrullDim_mathlib,
        lem:chart_ringKrullDim_le_one}
  Let \(\bar k\) be an algebraically closed field and let \(C : \Over(\Spec \bar k)\)
  be a curve over \(\bar k\) whose structure morphism is smooth of relative dimension
  one and locally of finite type, with \(C.\mathrm{left}\) integral.  Then
  \[
    \krullDim C.\mathrm{left} \;\le\; 1 .
  \]
\end{theorem}

\begin{proof}
  \uses{lem:krullDim_eq_iSup_coheight_mathlib, lem:ringKrullDim_stalk_eq_coheight,
        lem:stalk_standardSmooth_localization_of_smooth,
        lem:ringKrullDim_eq_height_mathlib, lem:height_le_ringKrullDim_mathlib,
        lem:chart_ringKrullDim_le_one}
        \leanok
  By \cref{lem:krullDim_eq_iSup_coheight_mathlib} the Krull dimension of the space
  \(C.\mathrm{left}\) is the supremum of the coheights of its points, so it suffices to
  prove \(\coheight(z) \le 1\) for every point \(z : C.\mathrm{left}\).  By the
  coheight bridge \cref{lem:ringKrullDim_stalk_eq_coheight} this coheight equals the
  Krull dimension of the stalk, so the goal becomes \(\krullDim \mathcal O_{C,z} \le
  1\).

  Fix \(z\).  By
  \cref{lem:stalk_standardSmooth_localization_of_smooth} the point \(z\) lies on an
  affine open chart \(V \ni z\) with coordinate ring \(S = \Gamma(C, V)\), a
  standard-smooth \(\bar k\)-algebra of relative dimension one, such that the stalk
  \(\mathcal O_{C,z}\) is the localisation of \(S\) at the prime \(\mathfrak p_z\)
  corresponding to \(z\).  Since \(C.\mathrm{left}\) is integral and \(V\) is a
  nonempty open, the section ring \(S\) is an integral domain.

  Now \cref{lem:ringKrullDim_eq_height_mathlib} identifies the stalk dimension with the
  height of the prime, \(\krullDim \mathcal O_{C,z} = \height \mathfrak p_z\), and
  \cref{lem:height_le_ringKrullDim_mathlib} bounds this height by the dimension of the
  chart ring, \(\height \mathfrak p_z \le \krullDim S\).  Finally
  \cref{lem:chart_ringKrullDim_le_one} gives \(\krullDim S \le 1\).  Chaining the
  inequalities,
  \[
    \krullDim \mathcal O_{C,z}
      = \height \mathfrak p_z
      \le \krullDim S
      \le 1 ,
  \]
  which is the required pointwise bound.  Taking the supremum over \(z\) yields
  \(\krullDim C.\mathrm{left} \le 1\).
\end{proof}
